\subsection{可作卷积的测度与函数的例子}

在命题 3 与命题 4 中，\(\mathcal{C}'(\mathrm{G})\) 与 \(\mathcal{M}(\mathrm{G})\) 分别赋予在 \(\mathcal{C}(\mathrm{G})\) 与 \(\mathcal{K}(\mathrm{G})\) 中的紧收敛拓扑。

命题 3. --- 假设 \(\chi\) 连续。设 \(\mu \in \mathcal{C}'(\mathbf{G})\)，\(f \in \mathcal{C}(\mathbf{X})\)。则：

\begin{enumerate}
\def\labelenumi{(\roman{enumi})}
\item
  \(\mu\) 与 \(f\) 关于 \(\beta\) 可作卷积。
\item
  No.~1 的公式 (3) 对每个 \(x \in \mathbf{X}\) 定义了一个连续的卷积 \(\mu *^{\beta}f\)，它正是元素 \(\pmb{\gamma}_{\chi}(\mu)f\)，即由连续表示 \(\pmb{\gamma}_{\chi}\)（\(\mathbf{G}\) 在 \(\mathcal{C}(\mathbf{X})\) 中）所定义者；此外，映射 \((\mu, f) \mapsto \mu *^{\beta}f\) 相对于 \(\mathcal{C}'(\mathbf{G})\) 的等度连续子集与 \(\mathcal{C}(\mathbf{X})\) 的紧子集是亚连续的。
\item
  若再有 \(f \in \mathcal{K}(\mathrm{X})\)，则 (ii) 中的乘积 \(\mu *^{\beta}f\) 属于 \(\mathcal{K}(\mathrm{X})\)，且映射 \((\mu, f) \mapsto \mu *^{\beta}f\) 相对于 \(\mathcal{C}'(\mathrm{G})\) 的等度连续子集与 \(\mathcal{K}(\mathrm{X})\) 的紧子集是亚连续的。
\end{enumerate}

我们知道 \(\mu\) 与 \(f\) 可作卷积（§3, No.~2 命题 8 (i)）。另一方面，采用 §2 的记号，我们有

\[
\boldsymbol {\gamma} _ {\chi} (\mu) f = \int \left(\boldsymbol {\gamma} _ {\chi} (s) f\right) d \mu (s) \in \mathcal {C} (\mathrm{X})
\]

因为 \(\mathcal{C}(\mathbf{X})\) 是完备的。特别地，对每个 \(x\in \mathbf{X}\)，

\[
\left(\boldsymbol {\gamma} _ {\chi} (\mu) f\right) (x) = \int \left(\boldsymbol {\gamma} _ {\chi} (s) f\right) (x) d \mu (s).
\]

把这一点与命题 2 (i) 以及 §2, No.~6 结合起来，便证明了 (ii)。最后，若 \(f \in \mathcal{K}(\mathrm{X})\)，则 \(\mu * (f \cdot \beta)\) 具有紧支集（§3, No.~2 命题 9），因此 \(\mu *^{\beta} f \in \mathcal{K}(\mathrm{X})\)。为此，考虑连续表示 \(U\)（\(\mathbf{G}\) 在完备化空间 \(\mathcal{K}(\mathrm{X})^{\wedge}\) 中的），它是由连续算子 \(\pmb{\gamma}_{\chi}(s)\)（在 \(\mathcal{K}(\mathbf{X})\) 中）连续延拓而得到的（§2, No.~1 备注 3）。设 S 为 \(\mu\) 的支集。诸函数 \(\pmb{\gamma}_{\chi}(s)f\)（\(s \in \mathbf{S}\)）的支集都包含在一个固定的紧集 K 中。集合 \(\mathcal{K}(\mathbf{X}, \mathbf{K})\) 是 \(\mathcal{K}(\mathbf{X})\) 的完备线性子空间。因此 \(U(\mu)f \in \mathcal{K}(\mathbf{X})\)。如前可见 \(U(\mu)f = \mu^{*\beta}f\)，而 (iii) 仍由 §2, No.~6 得出。

命题 4. --- 假设 G 在 X 中真作用且 \(\chi\) 连续。设 \(\mu \in \mathcal{M}(\mathrm{G})\)，\(f \in \mathcal{K}(\mathrm{X})\)。

\begin{enumerate}
\def\labelenumi{(\roman{enumi})}
\item
  \(\mu\) 与 \(f\) 关于 \(\beta\) 可作卷积。
\item
  No.~1 的公式 (3) 对每个 \(x \in \mathbf{X}\) 定义了一个连续的卷积 \(\mu *^{\beta}f\)。
\item
  映射 \((\mu, f) \mapsto \mu *^{\beta}f\)（从 \(\mathcal{M}(\mathrm{G}) \times \mathcal{K}(\mathrm{X})\) 到 \(\mathcal{C}(\mathrm{X})\) 中）相对于 \(\mathcal{M}(\mathrm{G})\) 的有界子集与 \(\mathcal{K}(\mathrm{X})\) 中包含在某个子空间 \(\mathcal{K}(\mathrm{X}, \mathrm{L})\) 内的紧子集（其中 \(\mathrm{L}\) 是 \(\mathrm{X}\) 的变动紧子集）是亚连续的。
\end{enumerate}

我们知道 \(\mu\) 与 \(f\) 可作卷积（§3, No.~2 命题 8 (ii)），且显然 (3) 中出现的积分对每个 \(x \in \mathbf{X}\) 都存在。设 \(\mathbf{K}\) 与 \(\mathbf{L}\) 是 \(\mathbf{X}\) 的两个紧子集。存在 \(\mathbf{H}\)（\(\mathbf{G}\) 的一个紧子集），使得关系 \(x \in \mathbf{K}\) 与 \(s^{-1}x \in \mathbf{L}\) 蕴含 \(s \in \mathbf{H}\)；取 \(\varphi \in \mathcal{H}_{+}(\mathbf{G})\) 满足 \(\varphi(s) = 1\)（对 \(s \in \mathbf{H}\)）。于是，对 \(f \in \mathcal{H}(\mathbf{X}, \mathbf{L})\) 与 \(x \in \mathbf{K}\)，

\[
\begin{array}{c} \int f (s ^ {- 1} x) \chi (s ^ {- 1}, x) d \mu (s) = \int f (s ^ {- 1} x) \chi (s ^ {- 1}, x) \varphi (s) d \mu (s) \\ = ((\varphi \cdot \mu) * ^ {\beta} f) (x). \end{array}
\]

因而 \(\int f(s^{-1}x)\chi(s^{-1},x)d\mu(s)\) 是 \(x\) 的连续函数，并定义了一个卷积 \(\mu * ^ {\beta} f\in \mathcal{C}(\mathrm{X})\)。此外，映射 \(\mu \mapsto \varphi \cdot \mu\)（从 \(\mathcal{M}(\mathrm{G})\) 到 \(\mathcal{C}'(\mathrm{G})\) 中）对于紧收敛拓扑是连续的。于是命题 3 (iii) 蕴含：映射 \((\mu ,f)\mapsto \mu * ^ {\beta} f\)（从 \(\mathcal{M}(\mathrm{G})\times \mathcal{K}(\mathrm{X},\mathrm{L})\) 到 \(\mathcal{C}(\mathrm{X})\) 中）对 X 的每个紧子集 L 相对于 \(\mathcal{K}(\mathrm{X},\mathrm{L})\) 的紧子集是亚连续的。特别地，映射 \((\mu ,f)\mapsto \mu * ^ {\beta} f\)（从 \(\mathcal{M}(\mathrm{G})\times \mathcal{K}(\mathrm{X})\) 到 \(\mathcal{C}(\mathrm{X})\) 中）是分别连续的。由于 \(\mathcal{K}(\mathrm{X})\) 是桶式空间，该映射相对于 \(\mathcal{M}(\mathrm{G})\) 的有界子集是亚连续的（TVS, III, §5, No.~3 命题 6）。

备注 1. --- 在命题 4 的假设下，映射 \(\mu \mapsto \mu *^{\beta}f\)（从 \(\mathcal{M}_{+}(\mathrm{G})\) 到 \(\mathcal{C}(\mathrm{X})\) 中）当 \(\mathcal{M}_{+}(\mathrm{G})\) 赋予淡拓扑时对每个 \(f\in \mathcal{K}(\mathrm{X})\) 是连续的。因为，设 K 为 X 的一个紧子集，S 为 \(f\) 的（紧）支集；由于 G 在 X 中真作用，由 \(s\in \mathrm{G}\) 中使得存在 \(x\in \mathrm{K}\) 满足 \(s^{-1}x\in \mathrm{S}\) 的那些 s 组成的集合是 G 的一个紧子集 L（GT, III, §4, No.~5 定理 1）。设 \(\varepsilon\) 为一个 \(>0\) 的数，\(\varphi\) 为 \(\mathcal{K}_{+}(\mathrm{G})\) 中在紧集 L 上等于 1 的函数，\(\mu_0\) 为 \(\mathcal{M}_{+}(\mathrm{G})\) 的一个元素；集合 \(\mathrm{W}_0\) 由测度 \(\mu \in \mathcal{M}_{+}(\mathrm{G})\) 组成，它们满足

\[
\left| \int \varphi (s) d \mu (s) - \int \varphi (s) d \mu_ {0} (s) \right| \leqslant \varepsilon
\]

是 \(\mu_0\) 在 \(\mathcal{M}_{+}(\mathrm{G})\) 中的一个邻域。另一方面，函数 \((s,x)\mapsto f(s^{-1}x)\chi (s^{-1},x)\) 在 \(\mathbf{L}\times \mathbf{K}\) 上一致连续，因此存在有限多个点 \(x_{i}\in \mathbf{K}\)（\(1\leqslant i\leqslant n\)），使得对每个 \(x\in \mathbf{K}\)，存在一个 \(i\)，使得

\[
\left| f (s ^ {- 1} x) \chi (s ^ {- 1}, x) - f (s ^ {- 1} x _ {i}) \chi (s ^ {- 1}, x _ {i}) \right| \leqslant \varepsilon
\]

对一切 \(s \in \mathbf{L}\) 成立。由于 \(\mu(\mathbf{L}) \leqslant \int \varphi(s) d\mu_0(s) + \varepsilon\) 对一切 \(\mu \in \mathbf{W}_0\) 成立，故还有

\[
\begin{array}{r l} \Big | \int f (s ^ {- 1} x) \chi (s ^ {- 1}, x) d \mu (s) - \int f (s ^ {- 1} x _ {i}) \chi (s ^ {- 1}, x _ {i}) d \mu (s) \Big | & \\ & \leqslant \varepsilon \Big (\int \varphi (s) d \mu_ {0} (s) + \varepsilon \Big) \end{array}
\]

对每个满足前一不等式的 x 以及每个 \(\mu \in W_{0}\) 成立。现在设 W 为 \(\mu_{0}\) 在 \(\mathcal{M}_{+}(G)\) 中的由测度 \(\mu \in W_{0}\) 组成的邻域，这些测度满足

\[
\left| \int f (s ^ {- 1} x _ {i}) \chi (s ^ {- 1}, x _ {i}) d \mu (s) - \int f (s ^ {- 1} x _ {i}) \chi (s ^ {- 1}, x _ {i}) d \mu_ {0} (s) \right| \leqslant \varepsilon
\]

对 \(1 \leqslant i \leqslant n\) 成立。显然，对每个测度 \(\mu \in \mathbf{W}\) 与每个 \(x \in \mathbf{K}\)，

\[
\begin{array}{r l} & {\Big | \int f (s ^ {- 1} x) \chi (s ^ {- 1}, x) d \mu (s) - \int f (s ^ {- 1} x) \chi (s ^ {- 1}, x) d \mu_ {0} (s) \Big |} \\ & {\qquad \leqslant \varepsilon \Big (2 \int \varphi (s) d \mu_ {0} (s) + 2 \varepsilon + 1 \Big),} \end{array}
\]

而由于 \(\varepsilon\) 是任意的，这就证明了我们的断言。

命题 5. --- 假设 \(\chi\) 是连续乘子且每个函数 \(\chi(s,\cdot)\) 有界。

\begin{enumerate}
\def\labelenumi{(\roman{enumi})}
\item
  函数 \(s \mapsto \rho(s) = \sup_{x \in X} \chi(s^{-1}, x)\)（定义在 \(G\) 上）是 \(> 0\) 的下半连续函数，且满足 \(\rho(st) \leqslant \rho(s)\rho(t)\)（对一切 \(s, t\)，\(G\) 中的）。
\item
  设 \(\mu \in \mathcal{M}^{\rho}(\mathrm{G})\) 且 \(f \in \mathrm{L}^{\infty}(\mathrm{X},\beta)\)。则 \(\mu\) 与 \(f\) 可作卷积，且 \(\mu *^{\beta} f\) 在局部几乎处处由 No.~1 的公式 (3) 给出。我们有 \(\mu *^{\beta} f \in \mathrm{L}^{\infty}(\mathrm{X},\beta)\)，且 \(\| \mu *^{\beta} f \|_{\infty} \leqslant \| \mu \|_{\rho} \| f \|_{\infty}\)。
\item
  若再有 \(f \in \mathcal{C}^{\infty}(\mathrm{X})\)（相应地，\(\overline{\mathcal{K}(\mathrm{X})}\)），则 No.~1 的公式 (3) 对每个 \(x\) 定义了一个卷积 \(\mu *^{\beta}f\)，它属于 \(\mathcal{C}^{\infty}(\mathrm{X})\)（相应地，\(\overline{\mathcal{K}(\mathrm{X})}\)）。
\item
  若 \(f \in \overline{\mathcal{K}(\mathrm{X})}\)，则由 (3) 定义的卷积 \(\mu *^{\beta}f\) 正是元素 \(\gamma_{\chi}(\mu)f\)，它由连续表示 \(\gamma_{\chi}\)（\(G\) 在 \(\overline{\mathcal{K}(\mathrm{X})}\) 中的）所定义。
\end{enumerate}

恒等式 \(\chi(st,x) = \chi(s,tx)\chi(t,x)\) 立即蕴含 \(\rho(st) \leqslant \rho(s)\rho(t)\)。另一方面，\(\rho\) 作为连续函数的上包络，是下半连续的。

设 \(\mu \in \mathcal{M}^{\rho}(\mathrm{G})\)。由 No.~1 的命题 1，\(\mu\) 与 1 可作卷积；命题 2 (i) 表明 \((|\mu|^{*\beta}1)(x) \leqslant \int_{\mathrm{G}} \rho(s)d|\mu|(s)\) 在局部 \(\beta\)-几乎处处成立。因此，若 \(f\) 是 \(\beta\)-可测的且 \(|f| \leqslant 1\)，则 \(\mu\) 与 \(f\) 可作卷积，且 \(\mathrm{N}_{\infty}(\mu^{*\beta}f) \leqslant \int \rho(s)d|\mu|(s)\)。此外，\(\mu^{*\beta}f\) 在局部几乎处处由 No.~1 的公式 (3) 给出，因为 No.~1 命题 2 的条件 (iii) 得以满足。这就蕴含 (ii)。

设 \(f\) 连续且绝对值以 1 为界。显然 (3) 中出现的积分对一切 \(x \in X\) 存在。我们来证明它们连续地依赖于 \(x\)。可以假设 \(\mu \geqslant 0\)。设 \(x_0 \in X\) 与 \(\varepsilon > 0\)。设 \(K\) 为 \(G\) 的一个紧子集，使得 \(\int_{G - K} \rho(s) d\mu(s) \leqslant \varepsilon\)。存在一个邻域 \(V\)（\(x_0\) 在 \(X\) 中的），使得 \(x \in V\) 蕴含

\[
\left| f \left(s ^ {- 1} x\right) \chi \left(s ^ {- 1}, x\right) - f \left(s ^ {- 1} x _ {0}\right) \chi \left(s ^ {- 1}, x _ {0}\right) \right| \leqslant \varepsilon / \mu (\mathrm{K})
\]

对 \(s \in \mathbf{K}\) 成立。于是，对 \(x \in \mathbf{V}\)，

\[
\begin{array}{r l} & {\Big | \int f (s ^ {- 1} x) \chi (s ^ {- 1}, x) d \mu (s) - \int f (s ^ {- 1} x _ {0}) \chi (s ^ {- 1}, x _ {0}) d \mu (s) \Big |} \\ & {\qquad \leqslant 2 \int_ {\mathrm{G-K}} \rho (s) d \mu (s) + \int_ {\mathrm{K}} \frac {\varepsilon}{\mu (\mathrm{K})} d \mu (s) \leqslant 3 \varepsilon ,} \end{array}
\]

由此得我们的断言。假设此外 \(f \in \overline{\mathcal{H}(\mathrm{X})}\)。设 \(\mathbf{H}\) 为 \(\mathbf{X}\) 的一个紧子集，使得 \(|f(y)| \leqslant \varepsilon\) 对 \(y \notin \mathbf{H}\) 成立。设 \(x \notin \mathbf{KH}\)。则 \(s^{-1}x \notin \mathbf{H}\) 对 \(s \in \mathbf{K}\) 成立，因此

\[
\begin{array}{r l} \Big | \int_ {\mathrm{G}} f (s ^ {- 1} x) \chi (s ^ {- 1}, x) d \mu (s) \Big | & \leqslant \int_ {\mathrm{G} - \mathrm{K}} \rho (s) d \mu (s) + \int_ {\mathrm{K}} \varepsilon \rho (s) d \mu (s) \\ & \leqslant \varepsilon \Big (1 + \int_ {\mathrm{G}} \rho (s) d \mu (s) \Big), \end{array}
\]

这就完成了 (iii) 的证明。

最后，若 \(f \in \overline{\mathcal{K}(\mathbf{X})}\)，则由于 \(\varepsilon_x \in \mathcal{M}^1(\mathbf{X})\) 对一切 \(x \in \mathbf{X}\) 成立，我们有

\[
\left(\boldsymbol {\gamma} _ {\chi} (\mu) f\right) (x) = \int \left(\boldsymbol {\gamma} _ {\chi} (s) f\right) (x) d \mu (s),
\]

从而 \(\gamma_{\chi}(\mu)f\) 就是由 (3) 定义的卷积 \(\mu *^{\beta}f\)。

命题 6. --- 假设 \(\chi\) 是连续乘子且每个函数 \(\chi(s, \cdot)\) 有界。设 \(\rho(s) = \sup_{x \in X} \chi(s^{-1}, x)\)。设 \(p\) 与 \(q\) 为一对共轭指数（\(1 \leqslant p < +\infty\)）。设 \(\mu \in \mathcal{M}^{\rho^{1/q}}(G)\) 且 \(f \in L^{p}(X, \beta)\)。则：

\begin{enumerate}
\def\labelenumi{(\roman{enumi})}
\item
  \(\mu\) 与 \(f\) 可作卷积；
\item
  卷积 \(\mu * ^ {\beta} f\) 在局部 \(\beta\)-几乎处处由公式 (3) 给出，且在局部 \(\beta\)-几乎处处等于一个函数 \(g\in \mathbf{L}^p (\mathbf{X},\beta)\)，使得 \(\| g\| _p\leqslant \| \mu \|_{\rho^{1 / q}}\| f\| _p;\) (iii) \(g\) 等于元素 \(\gamma_{\chi}(\mu)f\)，它由连续表示 \(\gamma_{\chi}\)（\(\mathbf{G}\) 在 \(\mathbf{L}^p (\mathbf{X},\beta)\) 中的）所定义。
\end{enumerate}

我们有

\[
\int^ {*} \| \gamma_ {\chi} (s) f \| _ {p} d | \mu | (s) \leqslant \left(\int^ {*} \rho (s) ^ {1 / q} d | \mu | (s)\right) \| f \| _ {p} <   + \infty
\]

（依据 §2, No.~5 的公式 (5)）。另一方面，映射 \(s \mapsto \pmb{\gamma}_{\chi}(s)f\)（从 G 到 \(\mathrm{L}^p (\mathbf{X},\beta)\) 中）是连续的（§2, No.~5 命题 9）。因此该映射是 \(\mu\)-可积的。设

\[
g = \int_ {\mathrm{G}} \left(\boldsymbol {\gamma} _ {\chi} (s) f\right) d \mu (s) \in \mathrm{L} ^ {p} (\mathrm{X}, \beta).
\]

我们有 \(\| g\| _p\leqslant \left(\int \rho^{1 / q}(s)d|\mu |(s)\right)\| f\| _p\)。把前面的备注应用于 \(|f|\)，可以看出映射 \(s\mapsto \varepsilon_s*\mid f|\)（从 G 到 \(\mathbf{L}^p (\mathbf{X},\beta)\) 中）是 \(\mu\)-可积的，因此对每个 \(h\in \mathcal{K}(\mathbf{X})\)，映射 \(s\mapsto \langle h,\varepsilon_s*(|f|\cdot \beta)\rangle\) 是 \(\mu\)-可积的。于是 §1, No.~5 的命题 7 证明 \(\mu\) 与 \(f\cdot \beta\) 可作卷积。此外，

\[
\begin{array}{r l} \int_ {\mathrm{X}} g (x) h (x) d \beta (x) & = \int_ {\mathrm{G}} d \mu (s) \int_ {\mathrm{X}} (\boldsymbol {\gamma} _ {\chi} (s) f) (x) h (x) d \beta (x) \\ & = \int_ {\mathrm{G}} \langle h, \varepsilon_ {s} * (f \cdot \beta) \rangle d \mu (s), \end{array}
\]

而由 §1, No.~5 的命题 7，最后这个积分等于 \(\langle h, \mu * (f \cdot \beta) \rangle\)。于是可见 g 是 \(\mu\) 与 f 的一个卷积。由命题 2 及其后的备注，这个卷积在局部 \(\beta\)-几乎处处由 (3) 给出。

按照记号的滥用，常常把陈述中的函数 g 之一记作 \(\mu * ^ {\beta} f\)，这使我们可以写出

\[
\left\| \mu * ^ {\beta} f \right\| _ {p} \leqslant \left\| \mu \right\| _ {\rho^ {1 / q}} \| f \| _ {p}.
\]

若 X 在无穷处可数，这种记法此外还是完全合理的。

推论. --- 在命题 6 的假设下，映射 \((\mu, f) \mapsto \mu *^{\beta} f\) 在 \(\mathrm{L}^{p}(\mathrm{X}, \beta)\) 上定义了 \(\mathcal{M}^{\rho^{1/q}}(\mathrm{G})\) 上的一个左模结构（\(1 \leqslant p \leqslant +\infty\)）。

这由命题 5、命题 6 以及卷积的结合性得出。

备注 2. --- 设 X 为局部紧空间，局部紧群 G 通过 \((x,s)\mapsto xs\) 连续地右作用于其上。设 \(\beta\) 为 X 上的正测度。设 \(\chi\) 为 \textgreater0 且定义在 \(G\times X\) 上的函数，它对 \(G\times X\) 上的每个测度可测，且 \(\boldsymbol{\delta}(s)\boldsymbol{\beta}=\chi(s,\cdot)\cdot\boldsymbol{\beta}\) 对每个 \(s\in G\) 成立。设 f 为 X 上局部 \(\beta\)-可积的函数，并设 \(\mu\) 为 G 上的测度。若 \(f\cdot\beta\) 与 \(\mu\) 可作卷积（对于映射 \((x,s)\mapsto xs\)，从 \(X\times G\) 到 X），则 \((f\cdot\beta)*\mu\) 以 \(\beta\) 为基。这时称 f 与 \(\mu\) 关于 \(\beta\) 可作卷积；\((f\cdot\beta)*\mu\) 关于 \(\beta\) 的每个密度都称为 f 与 \(\mu\) 关于 \(\beta\) 的一个卷积，记作 \(f * ^ {\beta} \mu\) 或简作 \(f*\mu\)。

设 \(G^{0}\) 为 G 的反群。通过 \((s,x)\mapsto xs\)，\(G^{0}\) 连续地左作用于 X。说 f 与 \(\mu\) 按前述意义可作卷积，等价于说 \(\mu\) 与 f 对于左作用于 X 的 \(G^{0}\) 可作卷积；而且卷积 \(f * ^ {\beta} \mu\) 不是别的，正是 \(\mu * ^ {\beta} f\) 对于左作用于 X 的 \(G^{0}\) 的卷积。另一方面，对 \(s\in G^{0}\) 有 \(\pmb{\gamma}(s)\beta = \chi(s^{-1},\cdot)\cdot \beta\)。于是 No.~1 与 No.~2 的结果可以立即翻译为关于乘积 \(f * ^ {\beta} \mu\) 的结果。特别地：

\begin{enumerate}
\def\labelenumi{\arabic{enumi})}
\tightlist
\item
  若 \(s \in \mathbf{G}\) 且 \(f\) 局部 \(\beta\)-可积，则 \(f\) 与 \(\varepsilon_s\) 可作卷积，且
\end{enumerate}

\[
(f * \varepsilon_ {s}) (x) = \chi (s ^ {- 1}, x) f (x s ^ {- 1}).\tag{4}
\]

\begin{enumerate}
\def\labelenumi{\arabic{enumi})}
\setcounter{enumi}{1}
\tightlist
\item
  若 \(f\) 与 \(\mu\) 可作卷积，且命题 2 的条件 (i)、(ii)、(iii) 之一成立，则 \(f *^{\beta} \mu\) 在局部 \(\beta\)-几乎处处由下式给出
\end{enumerate}

\[
(f * ^ {\beta} \mu) (x) = \int_ {G} f (x s ^ {- 1}) \chi (s ^ {- 1}, x) d \mu (s).\tag{5}
\]

我们把翻译其余陈述的任务留给读者。注意，若 \(\chi\) 连续且 \(\beta\) 的支集为 X，则

\[
\chi (t s, x) = \chi (s, x t) \chi (t, x) \quad (x \in \mathrm{X}; s, t \text {in} \mathrm{G}).\tag{6}
\]

